\documentclass[10pt,a4paper]{article}
\usepackage{../../styles/neet_style}

\begin{document}

% Title Header
\begin{center}
  {\Huge \textbf{NEET UG 2027: Master Practical Handbook}}\\[6pt]
  {\Large \textbf{Practical Chemistry: Qualitative Salt Analysis \& Volumetric Redox Titrations}}\\[4pt]
  {\small \textbf{Prepared, Compiled \& Authored by Satya Prakash} $\cdot$ Calibrated to Revised NMC Syllabus}
  \vspace{4pt}
  \hrule height 1.5pt
\end{center}

\pagestyle{fancy}
\fancyhf{}
\fancyhead[L]{\small \textbf{Practical Chemistry Lab Manual}}
\fancyhead[R]{\small Satya Prakash}
\fancyfoot[L]{\small \textbf{NEET UG 2027 Prep Series}}
\fancyfoot[C]{\thepage}
\fancyfoot[R]{\small \textbf{Chemistry Practical Notes}}

\section{Systematic Qualitative Analysis: Acidic Radicals (Anions)}

\begin{conceptbox}[Class I: Dilute $\mathrm{H_2SO_4}$ Soluble Anions ($\mathrm{CO_3^{2-}, S^{2-}, SO_3^{2-}, NO_2^-}$)]
Treated with dilute $\mathrm{H_2SO_4}$ to observe volatile evolved gases:
\begin{center}
\scriptsize
\begin{tabular}{|p{2.2cm}|p{1.2cm}|p{3.4cm}|p{7.2cm}|}
\hline
\textbf{Anion} & \textbf{Gas Evolved} & \textbf{Preliminary Observation} & \textbf{Confirmatory Test \& Governing Equation} \\
\hline
$\mathrm{CO_3^{2-}}$ (Carbonate) & $\mathrm{CO_2}$ & Brisk effervescence; odorless, colorless gas. & Turns lime water milky ($\mathrm{CaCO}_3 \downarrow$); prolonged bubbling clears milkiness forming soluble bicarbonate: $\mathrm{CaCO_3} + \mathrm{H_2O} + \mathrm{CO_2} \to \mathrm{Ca(HCO_3)_2}$. \\
\hline
$\mathrm{S^{2-}}$ (Sulfide) & $\mathrm{H_2S}$ & Colorless gas with rotten-egg smell. & Turns lead acetate paper black due to formation of lead sulfide ($\mathrm{PbS} \downarrow$); gives purple/violet coloration with \textbf{sodium nitroprusside}: $\mathrm{S^{2-}} + [\mathrm{Fe(CN)_5NO}]^{2-} \to [\mathrm{Fe(CN)_5NOS}]^{4-}$. \\
\hline
$\mathrm{SO_3^{2-}}$ (Sulfite) & $\mathrm{SO_2}$ & Colorless gas with suffocating odor of burning sulfur. & Turns acidified $\mathrm{K_2Cr_2O_7}$ paper green ($\mathrm{Cr^{3+}}$): $\mathrm{K_2Cr_2O_7} + \mathrm{H_2SO_4} + 3\mathrm{SO_2} \to \mathrm{Cr_2(SO_4)_3} + \mathrm{K_2SO_4} + \mathrm{H_2O}$. \\
\hline
$\mathrm{NO_2^-}$ (Nitrite) & $\mathrm{NO_2}$ & Pungent-smelling reddish brown fumes. & Starch-iodide paper turns deep blue: $2\mathrm{NO_2^-} + 2\mathrm{I^-} + 4\mathrm{H}^+ \to 2\mathrm{NO} + \mathrm{I_2} + 2\mathrm{H_2O}$; gives brown ring immediately with dilute $\mathrm{H_2SO_4}$ and $\mathrm{FeSO_4}$ (no conc. acid needed!). \\
\hline
\end{tabular}
\end{center}
\end{conceptbox}

\begin{conceptbox}[Class II: Concentrated $\mathrm{H_2SO_4}$ Anions ($\mathrm{Cl^-, Br^-, I^-, NO_3^-, CH_3COO^-}$)]
Treated with concentrated $\mathrm{H_2SO_4}$ and gently heated:
\begin{center}
\scriptsize
\begin{tabular}{|p{1.2cm}|p{3.8cm}|p{9.0cm}|}
\hline
\textbf{Anion} & \textbf{Observation} & \textbf{Confirmatory Reactions \& Diagnostic Identification} \\
\hline
$\mathrm{Cl^-}$ & Colorless pungent fumes ($\mathrm{HCl}$); dense white fumes with glass rod dipped in $\mathrm{NH_4OH}$ ($\mathrm{NH_4Cl}$). & \textbf{Chromyl Chloride Test:} Salt $+ \mathrm{K_2Cr_2O_7} + \mathrm{conc.\; H_2SO_4} \xrightarrow{\Delta}$ deep red vapors of chromyl chloride ($\mathrm{CrO_2Cl_2}$); pass into $\mathrm{NaOH} \to$ yellow solution of $\mathrm{Na_2CrO_4} \xrightarrow{\mathrm{Pb(CH_3COO)_2}}$ bright yellow precipitate of lead chromate ($\mathrm{PbCrO_4} \downarrow$). \\
\hline
$\mathrm{Br^-}$ & Reddish brown vapors ($\mathrm{Br_2}$) intensifying on adding $\mathrm{MnO_2}$. & \textbf{Organic Layer Test (Layer Test):} Salt solution $+ \text{chlorine water} + \mathrm{CHCl_3 / CCl_4} \xrightarrow{\text{shake}}$ organic layer turns \textbf{orange-brown} due to liberated free $\mathrm{Br_2}$ dissolved in non-polar layer ($2\mathrm{Br^-} + \mathrm{Cl_2} \to 2\mathrm{Cl^-} + \mathrm{Br_2}$). \\
\hline
$\mathrm{I^-}$ & Deep violet vapors ($\mathrm{I_2}$) turning starch paper blue. & \textbf{Layer Test:} Salt solution $+ \text{chlorine water} + \mathrm{CHCl_3} \xrightarrow{\text{shake}}$ organic layer turns \textbf{violet}. $\mathrm{AgNO_3}$ test gives bright yellow precipitate of $\mathrm{AgI}$ insoluble in concentrated $\mathrm{NH_4OH}$. \\
\hline
$\mathrm{NO_3^-}$ & Dense brown fumes of $\mathrm{NO_2}$ intensified by adding copper turnings. & \textbf{Brown Ring Test:} Salt solution $+ \text{freshly prepared } \mathrm{FeSO_4} \xrightarrow{\text{trickle conc. } \mathrm{H_2SO_4}} \mathbf{\text{Brown Ring at junction}}$: \newline $\mathrm{NO_3^-} + 3\mathrm{Fe^{2+}} + 4\mathrm{H}^+ \to \mathrm{NO} + 3\mathrm{Fe^{3+}} + 2\mathrm{H_2O}$, followed by: \newline $[\mathrm{Fe(H_2O)_6}]^{2+} + \mathrm{NO} \to \mathbf{[\mathrm{Fe(H_2O)_5(NO)}]^{2+}}$ (\textbf{Pentaaquanitrosyliron(I)}). \\
\hline
$\mathrm{CH_3COO^-}$ & Characteristic vinegar odor. & \textbf{Ester Test:} Salt $+ \text{ethanol} + \mathrm{conc.\; H_2SO_4} \xrightarrow{\Delta}$ sweet fruity smell of ethyl acetate ($\mathrm{CH_3COOC_2H_5}$). \textbf{Ferric Chloride Test:} Neutral $\mathrm{FeCl}_3 \to$ deep blood-red coloration ($\mathrm{[Fe_3(OH)_2(CH_3COO)_6]^+}$); boiling yields brown basic ferric acetate ppt. \\
\hline
\end{tabular}
\end{center}
\end{conceptbox}

\begin{formulabox}[Class III: Independent Group Anions ($\mathrm{SO_4^{2-}, PO_4^{3-}}$)]
Do not evolve gases with dilute or concentrated acids; identified by direct precipitation:
\begin{align*}
  \textbf{Sulfate } (\mathrm{SO_4^{2-}): } \quad & \text{Salt} + \mathrm{BaCl_2} \xrightarrow{\mathrm{dil.\; HCl}} \mathbf{\mathrm{BaSO_4} \downarrow} \quad (\text{White precipitate insoluble in conc. acid})\\
  \textbf{Phosphate } (\mathrm{PO_4^{3-}): } \quad & \text{Salt} + \mathrm{conc.\; HNO_3} + (\mathrm{NH_4})_2\mathrm{MoO_4} \xrightarrow{\Delta} \mathbf{(\mathrm{NH_4})_3[\mathrm{P(Mo_{3}O_{10})_4}] \downarrow} \quad (\textbf{Canary Yellow})
\end{align*}
\end{formulabox}

\section{Systematic Qualitative Analysis: Basic Radicals (Cations)}

\begin{conceptbox}[Systematic Cation Group Separation Scheme \& Group Reagents]
\begin{center}
\scriptsize
\begin{tabular}{|p{1.2cm}|p{2.2cm}|p{3.2cm}|p{2.6cm}|p{5.1cm}|}
\hline
\textbf{Group} & \textbf{Cations} & \textbf{Group Reagent} & \textbf{Precipitated As} & \textbf{Underlying Solubility Product ($K_{sp}$) Principle} \\
\hline
\textbf{Group 0} & $\mathrm{NH_4^+}$ & $\mathrm{NaOH}$ (heat) & Gas ($\mathrm{NH_3}$) & Identified by Nessler's reagent: brown ppt of iodide of Millon's base ($\mathrm{H_2N-Hg-O-Hg-I}$). \\
\hline
\textbf{Group I} & $\mathrm{Pb^{2+}}$ & Dilute $\mathrm{HCl}$ & Chlorides ($\mathrm{PbCl_2} \downarrow$) & Lowest $K_{sp}$ of chlorides; $\mathrm{PbCl_2}$ dissolves in hot water, recrystallizes on cooling. \\
\hline
\textbf{Group II} & $\mathrm{Pb^{2+}, Cu^{2+}, As^{3+}}$ & $\mathrm{H_2S}$ gas $+$ dilute $\mathrm{HCl}$ & Sulfides ($\mathrm{PbS, CuS, As_2S_3} \downarrow$) & \textbf{Common Ion Effect:} $\mathrm{H}^+$ suppresses $\mathrm{H_2S}$ dissociation, keeping $[\mathrm{S}^{2-}]$ very low; only Group II sulfides precipitate. \\
\hline
\textbf{Group III} & $\mathrm{Fe^{3+}, Al^{3+}}$ & $\mathrm{NH_4OH} + \mathbf{\mathrm{NH_4Cl}}$ & Hydroxides ($\mathrm{Fe(OH)_3, Al(OH)_3} \downarrow$) & \textbf{Common Ion Effect:} $\mathrm{NH_4^+}$ suppresses $\mathrm{NH_4OH}$ dissociation, keeping $[\mathrm{OH}^-]$ low to precipitate only Group III hydroxides. \\
\hline
\textbf{Group IV} & $\mathrm{Co^{2+}, Ni^{2+}},$ \newline $\mathrm{Mn^{2+}, Zn^{2+}}$ & $\mathrm{H_2S} + \mathbf{\mathrm{NH_4OH}}$ & Sulfides \newline ($\mathrm{CoS, NiS, MnS, ZnS} \downarrow$) & $\mathrm{OH}^-$ consumes $\mathrm{H}^+ \implies$ increases $[\mathrm{S}^{2-}]$, exceeding higher $K_{sp}$ of Group IV sulfides. \\
\hline
\textbf{Group V} & $\mathrm{Ba^{2+}, Sr^{2+}, Ca^{2+}}$ & $(\mathrm{NH_4})_2\mathrm{CO_3} +$ \newline $\mathrm{NH_4OH} + \mathrm{NH_4Cl}$ & Carbonates \newline ($\mathrm{BaCO_3, SrCO_3, CaCO_3} \downarrow$) & Carbonates precipitated in alkaline medium; $\mathrm{Mg}^{2+}$ remains in solution due to $\mathrm{NH_4Cl}$. \\
\hline
\textbf{Group VI} & $\mathrm{Mg^{2+}}$ & $\mathrm{Na_2HPO_4} + \mathrm{NH_4OH}$ & $\mathrm{Mg(NH_4)PO_4} \downarrow$ & White crystalline precipitate formed in alkaline ammoniacal medium. \\
\hline
\end{tabular}
\end{center}
\end{conceptbox}

\begin{formulabox}[High-Yield Confirmatory Tests \& Diagnostic Colorations]
\textbf{1. Nickel Ion ($\mathrm{Ni^{2+}}$) with Dimethylglyoxime (DMG):}
Rosy-red precipitate stabilized by intramolecular hydrogen bonding:
\begin{equation}
  \mathrm{Ni^{2+}} + 2\,\mathrm{DMG} + 2\,\mathrm{NH_4OH} \longrightarrow \mathbf{[\mathrm{Ni(DMG)_2}] \downarrow} \ (\text{Rosy Red}) + 2\,\mathrm{NH_4^+} + 2\,\mathrm{H_2O}
\end{equation}
\textbf{2. Ferric Ion ($\mathrm{Fe^{3+}}$) Confirmatory Tests:}
\begin{align*}
  \text{With Potassium Ferrocyanide: } & 4\mathrm{Fe^{3+}} + 3[\mathrm{Fe(CN)_6}]^{4-} \longrightarrow \mathbf{\mathrm{Fe_4}[\mathrm{Fe(CN)_6}]_3 \downarrow} \ (\textbf{Prussian Blue})\\
  \text{With Potassium Thiocyanate: } & \mathrm{Fe^{3+}} + \mathrm{SCN^-} \longrightarrow \mathbf{[\mathrm{Fe(SCN)}]^{2+}} \ (\textbf{Blood Red coloration})
\end{align*}
\textbf{3. Copper Ion ($\mathrm{Cu^{2+}}$) Confirmatory Tests:}
\begin{align*}
  \text{With Excess } \mathrm{NH_4OH}: \quad & \mathrm{Cu^{2+}} + 4\mathrm{NH_3} \longrightarrow \mathbf{[\mathrm{Cu(NH_3)_4}]^{2+}} \quad (\textbf{Azure Blue})\\
  \text{With Potassium Ferrocyanide: } & 2\mathrm{Cu^{2+}} + [\mathrm{Fe(CN)_6}]^{4-} \longrightarrow \mathbf{\mathrm{Cu_2}[\mathrm{Fe(CN)_6}] \downarrow} \quad (\textbf{Chocolate Brown})
\end{align*}
\textbf{4. Flame Test Colors of Alkaline Earth Metals:}
\begin{equation}
  \mathbf{\mathrm{Ba^{2+}} \implies \text{Apple Green}} \qquad \mathbf{\mathrm{Sr^{2+}} \implies \text{Crimson Red}} \qquad \mathbf{\mathrm{Ca^{2+}} \implies \text{Brick Red}}
\end{equation}
\end{formulabox}

\section{Volumetric Quantitative Redox Titrations}

\begin{formulabox}[$\mathrm{KMnO_4}$ vs Mohr's Salt Titration (Standardization)]
Mohr's salt is a primary standard double salt: $\mathbf{(\mathrm{NH_4})_2\mathrm{SO_4}\cdot\mathrm{FeSO_4}\cdot 6\mathrm{H_2O}}$ ($\mathrm{M.W.} = 392\,\mathrm{g/mol}$).
\begin{align*}
  \text{Oxidation Half: } & 5[\mathrm{Fe^{2+}} \longrightarrow \mathrm{Fe^{3+}} + e^-]\\
  \text{Reduction Half: } & \mathrm{MnO_4^-} + 8\mathrm{H}^+ + 5e^- \longrightarrow \mathrm{Mn^{2+}} + 4\mathrm{H_2O}\\
  \textbf{Overall Net Redox: } & \mathbf{\mathrm{MnO_4^-} + 5\mathrm{Fe^{2+}} + 8\mathrm{H}^+ \longrightarrow \mathrm{Mn^{2+}} + 5\mathrm{Fe^{3+}} + 4\mathrm{H_2O}}
\end{align*}
\begin{equation}
  \mathbf{\frac{M_1 V_1}{1} (\mathrm{KMnO_4}) = \frac{M_2 V_2}{5} (\text{Mohr's Salt})} \implies \mathbf{N_1 V_1 = N_2 V_2}
\end{equation}
\centerline{\small (Valence factor: $n = 5$ for $\mathrm{KMnO_4}$, $n = 1$ for $\mathrm{Fe^{2+}}$)}

\textbf{Crucial Practical Precautions:}
\begin{itemize}
  \item \textbf{Acidifying Agent:} Strictly use \textbf{dilute $\mathrm{H_2SO_4}$}.
  \begin{itemize}
    \item \textbf{Why NOT $\mathrm{HCl}$?} Permanganate oxidizes $\mathrm{HCl}$ to chlorine gas ($2\mathrm{MnO_4^-} + 16\mathrm{HCl} \to 2\mathrm{Mn^{2+}} + 5\mathrm{Cl_2} \uparrow + 8\mathrm{H_2O} + 6\mathrm{Cl^-}$), artificially consuming excess $\mathrm{KMnO_4}$!
    \item \textbf{Why NOT $\mathrm{HNO}_3$?} Nitric acid is a strong oxidizing agent itself and competes with $\mathrm{KMnO_4}$ in oxidizing $\mathrm{Fe}^{2+}$.
  \end{itemize}
  \item \textbf{Indicator:} $\mathrm{KMnO_4}$ acts as a \textbf{self-indicator}. The endpoint is marked by the appearance of a \textbf{permanent faint pink color} (excess single drop of $\mathrm{MnO_4^-}$).
\end{itemize}
\end{formulabox}

\begin{formulabox}[$\mathrm{KMnO_4}$ vs Oxalic Acid Titration]
Oxalic acid is a primary standard dihydrate: $\mathbf{(\mathrm{COOH})_2\cdot 2\mathrm{H_2O}}$ ($\mathrm{M.W.} = 126\,\mathrm{g/mol}, \text{ Equivalent Weight} = 63\,\mathrm{g/eq}$).
\begin{align*}
  \text{Oxidation Half: } & 5[\mathrm{C_2O_4^{2-}} \longrightarrow 2\mathrm{CO_2} \uparrow + 2e^-]\\
  \text{Reduction Half: } & 2[\mathrm{MnO_4^-} + 8\mathrm{H}^+ + 5e^- \longrightarrow \mathrm{Mn^{2+}} + 4\mathrm{H_2O}]\\
  \textbf{Overall Net Redox: } & \mathbf{2\mathrm{MnO_4^-} + 5\mathrm{C_2O_4^{2-}} + 16\mathrm{H}^+ \longrightarrow 2\mathrm{Mn^{2+}} + 10\mathrm{CO_2} \uparrow + 8\mathrm{H_2O}}
\end{align*}
\begin{equation}
  \mathbf{\frac{M_1 V_1}{2} (\mathrm{KMnO_4}) = \frac{M_2 V_2}{5} (\text{Oxalic Acid})} \implies \mathbf{N_1 V_1 = N_2 V_2}
\end{equation}
\centerline{\small (Valence factor: $n = 5$ for $\mathrm{KMnO_4}$, $n = 2$ for $\mathrm{C_2O_4^{2-}}$)}
\textbf{Crucial Practical Precautions:}
\begin{itemize}
  \item \textbf{Pre-Heating Requirement:} Oxalic acid solution must be heated to \textbf{$60-70^\circ\mathrm{C}$} before starting titration. At room temperature, the reaction is extremely slow.
  \item \textbf{Autocatalysis:} The initial disappearance of pink color is slow, but once the first few drops react, the reaction becomes rapid because product \textbf{$\mathrm{Mn}^{2+}$ acts as an autocatalyst}!
  \item \textbf{Do NOT boil ($>70^\circ\mathrm{C}$):} High temperature causes thermal decomposition of oxalic acid into $\mathrm{CO}_2, \mathrm{CO},$ and $\mathrm{H}_2\mathrm{O}$.
\end{itemize}
\end{formulabox}

\end{document}
